\chapter{Introduction}

\section{Prerequisites}

This text assumes familiarity with metric space theory and with linear algebra. 

\section{Linear algebra notation}

\begin{definition}
	Suppose $V$ is a finite dimensional vector space and $B = ( v_1, \dotsc, v_n )$ is a basis for $V$. Then any $v \in V$ can be written uniquely as $a_1v_1 + \dotsb + a_n v_n$ for some scalars $a_1, \dotsc, a_n \in \R$, and we define the \emph{representation of $v$ with respect to $B$}, denoted $[v]_B$, to be the column vector with $n$ entries that records the scalars $a_1, \dotsc, a_n$. In other words,
\[ [v]_B = \begin{bmatrix} a_1 \\ \vdots \\ a_n \end{bmatrix}. \]
\end{definition}

\begin{definition}
Suppose $\ell : V \to W$ is a linear map between finite dimensional vector spaces. Let $B = v_1, \dotsc, v_m$ be a basis for $V$ and let $C$ be a basis for $W$. The \emph{representation of $\ell$ with respect to $B$ and $C$}, denoted $[\ell]_{B,C}$, is the $n \times m$ matrix whose $i$th column of $[\ell]_{B,C}$ is $[\ell(v_i)]_C$. In other words,
\[ [\ell]_{B,C} = \begin{bmatrix} [\ell(v_1)]_C & \dotsb & [\ell(v_m)]_C \end{bmatrix}. \]
\end{definition}

\begin{lemma}
	If $\ell : V \to W$ is a linear map between finite dimensional vector spaces and $B$ and $C$ are bases for $V$ and $W$, respectively, then \[ [\ell]_{B,C}[v]_B = [\ell(v)]_C \]
	for any $v  \in V$. 
\end{lemma}

\begin{lemma}
	If $\ell : U \to V$ and $\ell' :  V \to W$ are linear maps between finite dimensional vector spaces, and $A, B$ and $C$ are bases for $U, V$, and $W$, respectively, then 
	\[ [\ell' \circ \ell]_{A,C} = [\ell']_{B,C}[\ell]_{A,B}. \]
\end{lemma}

\begin{lemma}
	If $\ell : V \to W$ is a linear map between finite dimensional vector spaces and $B$ and $C$ are bases for $V$ and $W$, respectively, then the rank of the matrix $[\ell]_{B,C}$ equals the dimension of the range $\ell(V)$.
\end{lemma}

\begin{definition}
	If $\ell : \R^m \to \R^n$ is a linear map, we just write $[\ell]$ to denote the matrix of $\ell$ with respect to the standard bases on $\R^m$ and $\R^n$.
\end{definition} 

\begin{lemma}
	If $A$ is a $n \times m$ matrix and $\ell_A : \R^m \to \R^n$ is the linear map $\ell_A(v) = Av$, then \[ [\ell_A] = A. \]
\end{lemma} 

\section{Little-oh notation}

Let $X$ be a metric space and let $x_0$ be a point in $X$. Also, let $g : X \to \R$ be a non-negative function; in other words, $g(x) \geq 0$ for all $x \in X$. \Cref{little-o-definition} below formulates what it means for another function $f : X \to \R$ to be ``little-oh of $g$ as $x$ tends to $x_0$.'' Intuitively, this condition means that $f$ is \emph{much smaller} than $g$ near $x_0$. 

\begin{definition} \label{little-o-definition}
	A function $f : X \to \R$ is \emph{little-oh of $g$ as $x$ tends to $x_0$}, written \[ f = o(g) \text{ as } x \to x_0 \] 
	if, for every $\epsilon > 0$, there exists an open neighborhood $U$ of $x_0$ such that $|f(x)| \leq \epsilon g(x)$ for all $x \in U$. When $x_0$ can be inferred from context, we write simply $f = o(g)$. 
\end{definition} 

\begin{remark}
	It's worth noting that the use of the symbol ``$=$'' above is mathematically abusive. The left-hand side of the ``$=$'' is a function and the right-hand side is a property of functions; of course, a function cannot literally be equal to a property. Rather, the ``$=$'' is being used here to mean that the function on the left-hand side has the property described on the right-hand side. As annoying as it is, this abuse of notation is fairly standard, so it's probably best to just get used to it. 
\end{remark}

\begin{lemma} \label{little-o-equivalent}
	Suppose that there exists a punctured neighborhood of $x_0$ in $X$ on which $g$ is nonzero. Prove that the following are equivalent conditions on a function $f : X \to \R$. 
	\begin{enumerate}[(a)]
		\item $f = o(g)$ as $x \to x_0$. 
		\item $\displaystyle \lim_{x \to x_0} \frac{f(x)}{g(x)} = 0$
	\end{enumerate}
\end{lemma}

\begin{exercise} \label{little-o-equivalent-exercise}
	Prove \cref{little-o-equivalent}.
\end{exercise}

\begin{exercise}
	Let $X = \R$. For each of the functions $f$ and $g$ described below, sketch graphs of $f$ and $g$ and then determine whether or not $f = o(g)$ as $x \to 0$.
	\begin{enumerate}[(a)]
		\item $f(x) = x$ and $g(x) = x^2$. 
		\item $f(x) = x$ and $g(x) = |x|$. 
		\item $f(x) = x^2$ and $g(x) = |x|$. 
	\end{enumerate}
\end{exercise}

\begin{exercise}
	Let $X = \R^2$. For each of the following functions $f$, determine whether or not $f = o(|x|)$ as $x \to 0$. 
	\begin{enumerate}[(a)]
		\item $f(x,y) = x - y$. 
		\item $f(x,y) = x^2 + y^2$. 
	\end{enumerate}
\end{exercise}

\begin{exercise} \label{little-o-vector-space}
	For a fixed non-negative function $g : X \to \R$, prove that the set of functions $f : X \to \R$ which are $o(g)$ as $x \to x_0$ is a vector space. In other words, verify the following three facts. 
	\begin{enumerate}[(1)]
		\item The zero function is $o(g)$. 
		\item If $c \in \R$ is a scalar and $f = o(g)$, then $cf = o(g)$. 
		\item If $f_1, f_2 = o(g)$, then $f_1 + f_2 = o(g)$. 
	\end{enumerate}
\end{exercise}

\begin{exercise} \label{little-o-multiplicative}
	Suppose $g_1, g_2 : X \to \R$ are non-negative functions and $f_1, f_2 : X \to \R$ are functions such that $f_1 = o(g_1)$ and $f_2 = o(g_2)$ as $x \to x_0$. Prove that $f_1 f_2 = o(g_1 g_2)$. 
\end{exercise}